\section{Propiedades de seguridad}
\label{sec:propiedades}
Nombres que usamos desde aqu\'i (ver Secci\'on~\ref{sec:protocolo}):
$N$ es el tama\~no del cat\'alogo, $i\in[1,N]$ el \'indice elegido,
$K_j$ la clave del registro $j$,
$\mathsf{ctx}$ el contexto inmutable que concatena de forma can\'onica
$\mathit{catalog\_id}$, $\mathit{version}$ y $\mathit{group\_id}$, y
$sid\in\{0,1\}^{128}$ el identificador \'unico de sesi\'on fresca.

El sistema debe satisfacer formalmente tres propiedades:

\begin{enumerate}
\item \textbf{Privacidad del receptor (selecci\'on oculta perfecta):}
Para cualesquiera dos \'indices $i_0, i_1 \in \{1,\dots,N\}$, la vista del
servidor es perfectamente indistinguible:
$\mathrm{Vista}_S(i_0) \equiv \mathrm{Vista}_S(i_1)$.
En nuestro protocolo, el mensaje $u = g^\alpha \cdot v^{-i}$ con $\alpha$
escogido uniformemente al azar en $\mathbb{Z}_q$ est\'a uniformemente
distribuido sobre el grupo $G$, enmascarando $i$ como una libreta de un solo
uso en el exponente. El servidor aprende exactamente cero informaci\'on sobre
$i$ en el sentido de la teor\'ia de la informaci\'on.

\item \textbf{Privacidad del emisor (transferencia \'unica bajo CDH):}
El cliente aprende $K_i$ y no obtiene informaci\'on computable sobre ninguna de
las dem\'as claves $\{K_{j\ne i}\}$.
Si un cliente malicioso pudiera extraer dos claves distintas en una sola
ejecuci\'on, podr\'ia calcular $v^\beta = g^{\beta^2}$, lo cual es
computacionalmente equivalente a resolver el problema Computational
Diffie-Hellman (CDH) en $G$ (seg\'un el Ejercicio~10.17 del libro
\cite{boneh-shoup-v06}), adem\'as de tener que romper el or\'aculo aleatorio
asociado a HKDF.

\item \textbf{Integridad, autenticidad y vinculaci\'on de contexto:}
Cualquier modificaci\'on no autorizada es detectada inmediatamente por el
cifrado autenticado (AEAD), arrojando la excepci\'on \code{InvalidTag}. Para
evitar ataques de confusi\'on de dominio o de traslaci\'on entre sesiones, se
definen dos datos asociados ($AAD$) diferenciados:
\begin{itemize}
\item \textbf{Dato asociado del cat\'alogo:}
$AAD_{\mathrm{cat}, j} = \mathsf{ctx}\Vert \mathrm{format\_u32}(j)$, utilizado
al cifrar cada registro $C_j$ de forma offline (no contiene $sid$ porque el
cat\'alogo es est\'atico e independiente de las sesiones).
\item \textbf{Dato asociado de la sesi\'on OT:}
$AAD_{\mathrm{ot}, j} = \mathsf{ctx}\Vert sid\Vert \mathrm{format\_u32}(j)$,
utilizado al encapsular cada clave $E_j$, ligando la respuesta a la sesi\'on
activa y al \'indice de posici\'on.
\end{itemize}
Asimismo, la huella $\mathrm{sha256}$ de la colecci\'on de textos cifrados
permite al cliente verificar que el cat\'alogo no haya sido alterado antes de
iniciar cualquier consulta.
\end{enumerate}
Las justificaciones matem\'aticas detalladas se presentan en la
Secci\'on~\ref{sec:protocolo}.
